\chapter{1950:Otto Diels and Kurt Alder}

\label{chap:chemistry1950}

The Nobel Prize in Chemistry for the year 1950 was awarded jointly to Otto Diels and Kurt Alder.

\textbf{Motivation:}

for their discovery and development of the diene synthesis

\section{Otto Diels}

\subsection*{Early Life and Education}

Otto Diels was born on 1876-01-23 in Hamburg, Germany.

\subsection*{Career and Major Research}

Diels studied chemistry at the University of Berlin under Emil Fischer, and in 1916 became professor and director of the Institute of Chemistry at the University of Kiel, where he remained until his retirement in 1945. With his student Kurt Alder he discovered the diene synthesis in 1928. Diels had earlier discovered the reactive ketenes (1905) and had isolated the product of a dehydrogenation of cholesterol that later proved to be a steroid.

\subsection*{Nobel Prize-winning Work}

The work for which Otto Diels was awarded the Nobel Prize in Chemistry in 1950 was described by the Nobel Committee as: "for their discovery and development of the diene synthesis".

The diene synthesis, now called the Diels-Alder reaction, joins a conjugated diene and a dienophile in a single step to form a six-membered ring. Reported by Diels and Alder in 1928, it provides a general route to cyclic organic compounds.

\subsection*{Significance and Impact}

The Diels-Alder reaction became one of the most important methods in synthetic organic chemistry for constructing six-membered rings. It was applied widely in the chemical industry, including the manufacture of synthetic rubber and plastics.

\subsection*{Later Career and Honors}

Diels retired from Kiel in 1945 and remained in Kiel until his death on 1954-03-07.

\subsection*{Legacy}

The Diels-Alder reaction, named for Diels and Alder, remains a standard and general reaction in organic synthesis, used in both academic research and industry.

\subsection*{Selected Publications}

\begin{itemize}

\item \emph{Synthesen in der hydroaromatischen Reihe}, Justus Liebigs Annalen der Chemie, 1928.

\end{itemize}

\clearpage

\section{Kurt Alder}

\subsection*{Early Life and Education}

Kurt Alder was born on 1902-07-10 in K{\"o}nigsh{\"u}tte, Germany (now Chorz{\'o}w, Poland).

\subsection*{Career and Major Research}

Alder studied chemistry in Berlin and took his doctorate at the University of Kiel in 1926 under Otto Diels. He worked with Diels on the diene synthesis, then joined the I. G. Farbenindustrie works at Leverkusen, where he worked on the synthetic rubber Buna. In 1940 he became professor of chemistry at the University of Cologne.

\subsection*{Nobel Prize-winning Work}

The work for which Kurt Alder was awarded the Nobel Prize in Chemistry in 1950 was described by the Nobel Committee as: "for their discovery and development of the diene synthesis".

Alder contributed to the systematic exploration of the diene synthesis, establishing its scope and its rules of orientation, and applying it to the construction of many ring systems.

\subsection*{Significance and Impact}

Alder's investigations extended the diene synthesis from a single observation into a general reaction of predictable outcome. The reaction proved essential in the synthesis of natural products and industrial polymers.

\subsection*{Later Career and Honors}

Alder remained at the University of Cologne as professor and director of its Institute of Chemistry. He died on 1958-06-20 in Cologne.

\subsection*{Legacy}

The Diels-Alder reaction and the Alder-ene reaction, both discovered in his laboratories, are named in his honor and remain widely used in organic synthesis.

\subsection*{Selected Publications}

\begin{itemize}

\item \emph{Synthesen in der hydroaromatischen Reihe}, Justus Liebigs Annalen der Chemie, 1928.

\end{itemize}

\clearpage

\section*{Chapter Summary}

The 1950 prize recognized the discovery and development of the diene synthesis. Otto Diels and Kurt Alder showed in 1928 that a conjugated diene and a dienophile combine in one step to form a six-membered ring, and the reaction became one of the most widely used methods in synthetic organic chemistry.
